\input eplain
\pdfpagewidth=148mm
\pdfpageheight=210mm
\hsize=128mm
\vsize=190mm

\pdfhorigin=0pt
\pdfvorigin=0pt
\hoffset=10mm
\voffset=10mm
\noindent 2 3 34 \leaders\hbox to 10pt {\hss.\hss}  \hfill \xref{node0}
\smallskip
\noindent 2 3 40 \leaders\hbox to 10pt {\hss.\hss}  \hfill \xref{node1}
\smallskip
\noindent 2 3 41 \leaders\hbox to 10pt {\hss.\hss}  \hfill \xref{node2}
\smallskip
\noindent 2 4 23 \leaders\hbox to 10pt {\hss.\hss}  \hfill \xref{node4}
\smallskip
\noindent 2 4 4 \leaders\hbox to 10pt {\hss.\hss}  \hfill \xref{node3}
\smallskip
\noindent 2 4 42 \leaders\hbox to 10pt {\hss.\hss}  \hfill \xref{node5}
\smallskip
\noindent 2 4 47 \leaders\hbox to 10pt {\hss.\hss}  \hfill \xref{node6}
\smallskip
\noindent 2 5 12 \leaders\hbox to 10pt {\hss.\hss}  \hfill \xref{node9}
\smallskip
\noindent 2 5 15 \leaders\hbox to 10pt {\hss.\hss}  \hfill \xref{node10}
\smallskip
\noindent 2 5 25 \leaders\hbox to 10pt {\hss.\hss}  \hfill \xref{node11}
\smallskip
\noindent 2 5 28 \leaders\hbox to 10pt {\hss.\hss}  \hfill \xref{node12}
\smallskip
\noindent 2 5 39 \leaders\hbox to 10pt {\hss.\hss}  \hfill \xref{node13}
\smallskip
\noindent 2 5 41 \leaders\hbox to 10pt {\hss.\hss}  \hfill \xref{node14}
\smallskip
\noindent 2 5 49 \leaders\hbox to 10pt {\hss.\hss}  \hfill \xref{node15}
\smallskip
\noindent 2 5 7 \leaders\hbox to 10pt {\hss.\hss}  \hfill \xref{node7}
\smallskip
\noindent 2 5 73 \leaders\hbox to 10pt {\hss.\hss}  \hfill \xref{node16}
\smallskip
\noindent 2 5 75 \leaders\hbox to 10pt {\hss.\hss}  \hfill \xref{node17}
\smallskip
\noindent 2 5 78 \leaders\hbox to 10pt {\hss.\hss}  \hfill \xref{node18}
\smallskip
\noindent 2 5 9 \leaders\hbox to 10pt {\hss.\hss}  \hfill \xref{node8}
\smallskip
\noindent 2 6 15 \leaders\hbox to 10pt {\hss.\hss}  \hfill \xref{node19}
\smallskip
\noindent 2 6 17 \leaders\hbox to 10pt {\hss.\hss}  \hfill \xref{node20}
\smallskip
\noindent 2 6 21 \leaders\hbox to 10pt {\hss.\hss}  \hfill \xref{node21}
\smallskip
\noindent 2 6 23 \leaders\hbox to 10pt {\hss.\hss}  \hfill \xref{node22}
\smallskip
\noindent 2 6 25 \leaders\hbox to 10pt {\hss.\hss}  \hfill \xref{node23}
\smallskip
\noindent 2 6 28 \leaders\hbox to 10pt {\hss.\hss}  \hfill \xref{node24}
\smallskip
\noindent 2 6 39 \leaders\hbox to 10pt {\hss.\hss}  \hfill \xref{node25}
\smallskip
\noindent 2 6 43 \leaders\hbox to 10pt {\hss.\hss}  \hfill \xref{node26}
\smallskip
\noindent 2 6 61 \leaders\hbox to 10pt {\hss.\hss}  \hfill \xref{node27}
\smallskip
\noindent 2 6 65 \leaders\hbox to 10pt {\hss.\hss}  \hfill \xref{node28}
\smallskip
\noindent 2 6 86 \leaders\hbox to 10pt {\hss.\hss}  \hfill \xref{node29}
\smallskip
\vfill \break
{\medskip\noindent\bf\font\titlefont=cmss17 at 20pt\titlefont Timeline}\medskip
{\noindent \bf 2026-02-09:} 2 3 34, 2 3 40, 2 3 41\smallskip
{\noindent \bf 2026-02-10:} 2 4 23, 2 4 4, 2 4 42, 2 4 47, 2 5 12, 2 5 15, 2 5 25, 2 5 28, 2 5 39, 2 5 41, 2 5 49, 2 5 7, 2 5 73, 2 5 75, 2 5 78, 2 5 9\smallskip
{\noindent \bf 2026-02-11:} 2 5 41, 2 5 49, 2 5 73, 2 5 75, 2 5 78, 2 6 15, 2 6 17, 2 6 21, 2 6 23, 2 6 25, 2 6 28, 2 6 39, 2 6 43, 2 6 61, 2 6 65, 2 6 86\smallskip
\vfill \break
\xrdef{node0}
{\medskip\noindent\bf\font\titlefont=cmss17 at 20pt\titlefont 2 3 34}\medskip
Find the following limit or state that it does not exist:
$$\
\lim_{x \to -5} {\
{x^2 + 7x + 10}
\over
{x + 5}
}
$$
{\medskip\noindent\bf\font\titlefont=cmitt12 at 17pt\titlefont 2026-02-09}\medskip
{\smallskip\noindent\bf\font\titlefont=cmr12 at 13pt\titlefont $\bullet$ 15:15: 2.3.34 no problem}\smallskip
This is a matter of factoring and canceling
\par

\xrdef{node1}
{\medskip\noindent\bf\font\titlefont=cmss17 at 20pt\titlefont 2 3 40}\medskip
Find the following limit or state that it does
not exist.
$$\
\lim_{h \to 0} {\
{{1 \over 2 + h} - {1 \over 2}} \over h
}
$$
{\medskip\noindent\bf\font\titlefont=cmitt12 at 17pt\titlefont 2026-02-09}\medskip
{\smallskip\noindent\bf\font\titlefont=cmr12 at 13pt\titlefont $\bullet$ 15:16: fractions in numerator, tricky}\smallskip
I keep doing this one wrong. This is the one that showed up on the quiz that I failed.
\par

My first attempt was 0, but that was wrong because I forgot about one of the terms.
\par

Huh. I am really not understanding the algebra for this
\par

okay great. trick I was missing was to distribute the denominator properly, and then properly combine fractions in the numerator.
\par

\xrdef{node2}
{\medskip\noindent\bf\font\titlefont=cmss17 at 20pt\titlefont 2 3 41}\medskip
Find the following limit or state that it does not exist.
$$\
\lim_{x \to 169} {\
{\sqrt{x} - 13}
\over
{x - 169}
}
$$
{\medskip\noindent\bf\font\titlefont=cmitt12 at 17pt\titlefont 2026-02-09}\medskip
{\smallskip\noindent\bf\font\titlefont=cmr12 at 13pt\titlefont $\bullet$ 15:35: 2.3.41 remember $13^2$ is 169}\smallskip
This can be factored and cancelled because the square root of 169 is 13.
\par

\xrdef{node3}
{\medskip\noindent\bf\font\titlefont=cmss17 at 20pt\titlefont 2 4 4}\medskip
Consider the function $F(x) = {f(x) \over g(x)}$ with $g(a) = 0$.
Does $F$ necessarily have a vertical asymptote at $x = a$?
Explain your reasoning.
{\medskip\noindent\bf\font\titlefont=cmitt12 at 17pt\titlefont 2026-02-10}\medskip
{\smallskip\noindent\bf\font\titlefont=cmr12 at 13pt\titlefont $\bullet$ 09:04: initial thoughts}\smallskip
In other words, if there is a zero in the denominator, does that necessarily have a vertical asymptote?
\par

if g(a) were a constant, there wouldn't be a asymptote right? Also, there could be expressions that cancel.
\par

{\smallskip\noindent\bf\font\titlefont=cmr12 at 13pt\titlefont $\bullet$ 10:08: a 0 in denominator doesn't necessarily have asymptote}\smallskip
As shown in example 5 in this chapter, you can have an undefined point that is not a vertical asymptote.
\par

\xrdef{node4}
{\medskip\noindent\bf\font\titlefont=cmss17 at 20pt\titlefont 2 4 23}\medskip
Determine the following limits, using $\infty$ or $-\infty$ when
appropriate, or state they do not exist.
\par a. $$\
\lim_{x \to 3^{+}}{
x - 6
\over
(x - 3)^2
}
$$
\par b. $$\
\lim_{x \to 3^{-}}{
x - 6
\over
(x - 3)^2
}
$$
\par c.$$\
\lim_{x \to 3}{
x - 6
\over
(x - 3)^2
}
$$
{\medskip\noindent\bf\font\titlefont=cmitt12 at 17pt\titlefont 2026-02-10}\medskip
{\smallskip\noindent\bf\font\titlefont=cmr12 at 13pt\titlefont $\bullet$ 09:07: initial thoughts}\smallskip
I plotted this one last night, and it looks like it approaches negative infinity.
\par

I am still tripped up on left/right sided limits. I guess we can say that both sides on the bottom approach zero from a positive angle?
\par

{\smallskip\noindent\bf\font\titlefont=cmr12 at 13pt\titlefont $\bullet$ 10:14: 2.4.23 initial work}\smallskip
Because the denominator is squared, it will approach 0 from positive values on both sides. The both sided limit exists. Becausee the numerator is negative, the value is negative infinity.
\par

\xrdef{node5}
{\medskip\noindent\bf\font\titlefont=cmss17 at 20pt\titlefont 2 4 42}\medskip
Determine the following limit, using $\infty$ or $-\infty$
when appropriate, or state that it does not exist.
$$\
\lim_{\theta \to {7\pi \over 2}^{+}}-{
1
\over
2
}\tan \theta
$$
{\medskip\noindent\bf\font\titlefont=cmitt12 at 17pt\titlefont 2026-02-10}\medskip
{\smallskip\noindent\bf\font\titlefont=cmr12 at 13pt\titlefont $\bullet$ 09:10: initial thoughts}\smallskip
I forget the behavior of tan. I think it's bounded between $\pi \over 2$ positive or negative.
\par

I think there theorems for this?
\par

{\smallskip\noindent\bf\font\titlefont=cmr12 at 13pt\titlefont $\bullet$ 10:40: 2.4.42 initial work}\smallskip
this is a trig limit for $tan$, for a big theta value. I came up with a potential solution by studying example 6 of 2.4. First, I rewrote the tan expression, then wrapped the bounds of the limit.
\par

The $\tan$ function can be rewritten as ${\sin \theta \over \cos \theta}$. sin evaluates to -1, and the cos evals to 0, plus the negative coefficient. at $7\pi \over 2$, cosine, approaches 0 positively from the RHS.
\par

I did some wrapping initially because my graphing calculator had floating point precision issues. lol.
\par

I want to say it is positive infinity, but I'm a little unsure how to distribute the negatives signs.
\par

\xrdef{node6}
{\medskip\noindent\bf\font\titlefont=cmss17 at 20pt\titlefont 2 4 47}\medskip
Find all vertical asymptotes, $x = a$, of the following
function. For each value $a$, evaluate
$$\lim_{x \to a^{+}}{f(x)}$$
$$\lim_{x \to a^{-}}{f(x)}$$ and
$$\lim_{x \to a}{f(x)}$$
Use $\infty$ or $-\infty$ when
appropriate.
$$\
f(x) = {
x^2 - 7x + 10
\over
x^2 - 6x + 8
}
$$
{\medskip\noindent\bf\font\titlefont=cmitt12 at 17pt\titlefont 2026-02-10}\medskip
{\smallskip\noindent\bf\font\titlefont=cmr12 at 13pt\titlefont $\bullet$ 09:14: initial thoughts}\smallskip
ah, is this the thing where you look at the highest order of each? there's a rule for polynomials.
\par

They want the lefthand and righthand sided limits.
\par

Oops, was there an $a$ value I needed to add?
\par

{\smallskip\noindent\bf\font\titlefont=cmr12 at 13pt\titlefont $\bullet$ 11:05: initial work}\smallskip
I think I forgot to supply $a$, but I could still factor and simplify to find asymptotes.
\par

I was able to cancel some terms, and apparently this means those zeros don't form asymptotes. Still don't understand why that works. I plotted this function and confirmed there is indeed only one VA.
\par

The limits go +/- infinity at that asymptote, so that limit doesn't exist (at 4).
\par

I'm not going to log on to find $a$, I'm just going to move onto doing 2.5.
\par

\xrdef{node7}
{\medskip\noindent\bf\font\titlefont=cmss17 at 20pt\titlefont 2 5 7}\medskip
Determine the following limit:
$$\
\lim_{t \to \infty}{
-13t^{-3}
}
$$
{\medskip\noindent\bf\font\titlefont=cmitt12 at 17pt\titlefont 2026-02-10}\medskip
{\smallskip\noindent\bf\font\titlefont=cmr12 at 13pt\titlefont $\bullet$ 09:16: initial thoughts}\smallskip
negative exponent means it becomes a denominator
\par

$t^3$ will approach infinity. but since it's in the denominator, it will approach 0.
\par

negation makes it negative infinity.
\par

{\smallskip\noindent\bf\font\titlefont=cmr12 at 13pt\titlefont $\bullet$ 18:08: 2.5.7 work}\smallskip
The main term is $t^{-3} = 1/t^3$, which goes to 0, making the whole thing 0.
\par

\xrdef{node8}
{\medskip\noindent\bf\font\titlefont=cmss17 at 20pt\titlefont 2 5 9}\medskip
Evaluate the following limit:
$$\
\lim_{x \to \infty}{
6 + {2 \over x^2}
}
$$
{\medskip\noindent\bf\font\titlefont=cmitt12 at 17pt\titlefont 2026-02-10}\medskip
{\smallskip\noindent\bf\font\titlefont=cmr12 at 13pt\titlefont $\bullet$ 11:10: initial thoughts}\smallskip
x goes to infinity, the square in the denominator goes to 0. The constant (6) remains. I think that's right?
\par

{\smallskip\noindent\bf\font\titlefont=cmr12 at 13pt\titlefont $\bullet$ 18:13: initial work}\smallskip
This one has a term under the numerator, making it go to zero, leaving the constant 6 remaining.
\par

\xrdef{node9}
{\medskip\noindent\bf\font\titlefont=cmss17 at 20pt\titlefont 2 5 12}\medskip
Determine the following limit at infinity:
$$\
\lim_{x \to \infty}{
7 + 7x + 7x^4
\over
x^4
}
$$
{\medskip\noindent\bf\font\titlefont=cmitt12 at 17pt\titlefont 2026-02-10}\medskip
{\smallskip\noindent\bf\font\titlefont=cmr12 at 13pt\titlefont $\bullet$ 11:12: initial thoughts}\smallskip
I think this is the one where you divide by the denominator? That trick?
\par

Weird that the powers are in reverse order. Why are they all 7s?
\par

I think you just consider the highest order powers and their coefficients. So, I think it approaches 7?
\par

{\smallskip\noindent\bf\font\titlefont=cmr12 at 13pt\titlefont $\bullet$ 18:16: initial work}\smallskip
This is a rational polynomial expression, so we do the trick where we divide by the highest polynomial in the denominator: $x^4$. this blows away everything but 7.
\par

\xrdef{node10}
{\medskip\noindent\bf\font\titlefont=cmss17 at 20pt\titlefont 2 5 15}\medskip
Suppose the function $g$ satisfies the inequality
$$\
30 - {11 \over x^4} \le 30 + {11 \over x^4}
$$
for all non-zero values of $x$. Evaluate
$$\lim_{x \to \infty}{g(x)}$$
and
$$\lim_{x \to -\infty}{g(x)}$$
{\medskip\noindent\bf\font\titlefont=cmitt12 at 17pt\titlefont 2026-02-10}\medskip
{\smallskip\noindent\bf\font\titlefont=cmr12 at 13pt\titlefont $\bullet$ 11:14: initial thoughts}\smallskip
Instead of an explicit equation, it gives an inequality. The function sits between two functions.
\par

Is this a squeeze theorem thing? You can solve the limits in between, and if they match there's a limit.
\par

eyeballing it: $x^4$ approaches infinity on both sides, which makes it approach 0 on both sides. Both approach 30.
\par

{\smallskip\noindent\bf\font\titlefont=cmr12 at 13pt\titlefont $\bullet$ 18:20: initial work}\smallskip
The hardest part of this one for me was understanding the inequality correctly. I'm still not 100% convinced, but I am pretty sure it is somewhere in between, so you can use the squeeze theorem. The term nlows away at both infinities, so it is just 30.
\par

\xrdef{node11}
{\medskip\noindent\bf\font\titlefont=cmss17 at 20pt\titlefont 2 5 25}\medskip
Determine the following limit:
$$\
\lim_{x \to \infty}{
15x^3 + 8x^2 - 8x
\over
21x^3 + 8x^2 + 4x + 5
}
$$
{\medskip\noindent\bf\font\titlefont=cmitt12 at 17pt\titlefont 2026-02-10}\medskip
{\smallskip\noindent\bf\font\titlefont=cmr12 at 13pt\titlefont $\bullet$ 11:19: initial thoughts}\smallskip
Is it only the highest order terms to consider? $15/21$ conveniently simplifies to $5/7$.
\par

{\smallskip\noindent\bf\font\titlefont=cmr12 at 13pt\titlefont $\bullet$ 18:30: initial work}\smallskip
This one I worked out pretty quickly, but hesitated on because I wasn't convinced it would go to a constant, and I was misrememnering theorem 2.7.
\par

I did it the "long" way by dividing all the terms, but plugging in theorem 2.7 would have been faster.
\par

\xrdef{node12}
{\medskip\noindent\bf\font\titlefont=cmss17 at 20pt\titlefont 2 5 28}\medskip
Determine the following limit:
$$\
\lim_{x \to \infty} {
3x^5 + 26
\over
4x^6 + 9x^5 - 5x^3
}
$$
{\medskip\noindent\bf\font\titlefont=cmitt12 at 17pt\titlefont 2026-02-10}\medskip
{\smallskip\noindent\bf\font\titlefont=cmr12 at 13pt\titlefont $\bullet$ 11:21: initial thoughts}\smallskip
These polynomial orders are different here. There's a theorem with a rule about. Intuitively, the denominator approaches infinity faster than the numerator. So does that mean it approaches 0?
\par

{\smallskip\noindent\bf\font\titlefont=cmr12 at 13pt\titlefont $\bullet$ 18:42: initial work}\smallskip
If I'm understanding theorem 2.7 correctly, then the answer to this one is simply zero.
\par

\xrdef{node13}
{\medskip\noindent\bf\font\titlefont=cmss17 at 20pt\titlefont 2 5 39}\medskip
Evalulate
$$\lim_{x \to \infty}{f(x)}$$
and
$$\lim_{x \to -\infty}{f(x)}$$
for the following function. Then give the
horizontal asymptote of f (if any).
$$\
f(x) = {
8x^2 - 9x + 7
\over
2x^2 + 3
}
$$
{\medskip\noindent\bf\font\titlefont=cmitt12 at 17pt\titlefont 2026-02-10}\medskip
{\smallskip\noindent\bf\font\titlefont=cmr12 at 13pt\titlefont $\bullet$ 11:55: initial thoughts}\smallskip
Is this again just about looking at the highest power? Looks like it would be 4.
\par

{\smallskip\noindent\bf\font\titlefont=cmr12 at 13pt\titlefont $\bullet$ 18:34: initial work}\smallskip
Via 2.7 theorem, these should evalulate to be 4, with a horizontal asymptote as y =4.
\par

\xrdef{node14}
{\medskip\noindent\bf\font\titlefont=cmss17 at 20pt\titlefont 2 5 41}\medskip
Determine
$$\lim_{x \to \infty}{f(x)}$$
and
$$\lim_{x \to -\infty}{f(x)}$$
for the following function. Then give the
horizontal asymptote of $f$ (if any).
$$\
f(x) = {
3x^5 - 2
\over
x^6 + 7x^4
}
$$
{\medskip\noindent\bf\font\titlefont=cmitt12 at 17pt\titlefont 2026-02-10}\medskip
{\smallskip\noindent\bf\font\titlefont=cmr12 at 13pt\titlefont $\bullet$ 11:56: initial thoughts}\smallskip
Powers differ by 1, with smaller (odd) power on top. Odd powers are negative as it approaches negative infinity, positive as it approaches positive infinity. Since it's asking to eval both sides. My guess is negative/positve for negative/positive infinity.
\par

As for the asymptotes: you can factor the denominator to get a zero that way maybe? looks like $\sqrt(7)$
\par

{\medskip\noindent\bf\font\titlefont=cmitt12 at 17pt\titlefont 2026-02-11}\medskip
{\smallskip\noindent\bf\font\titlefont=cmr12 at 13pt\titlefont $\bullet$ 11:01: initial work}\smallskip
I just plugged in 2.7 theorem. it goes to 0 because the numerator order is smaller.
\par

\xrdef{node15}
{\medskip\noindent\bf\font\titlefont=cmss17 at 20pt\titlefont 2 5 49}\medskip
Evaluate
$$\lim_{x \to \infty}{f(x)}$$
and
$$\lim_{x \to -\infty}{f(x)}$$
for the following function. Then give the
horizontal asymptote of $f$ (if any).
$$\
f(x) = {
\root 3 \of {x^6 + 64}
\over
3x^2 + \sqrt {5x^4 + 9}
}
$$
{\medskip\noindent\bf\font\titlefont=cmitt12 at 17pt\titlefont 2026-02-10}\medskip
{\smallskip\noindent\bf\font\titlefont=cmr12 at 13pt\titlefont $\bullet$ 12:00: initial thoughts}\smallskip
This one has roots, and I forget how to handle limits of those.
\par

{\medskip\noindent\bf\font\titlefont=cmitt12 at 17pt\titlefont 2026-02-11}\medskip
{\smallskip\noindent\bf\font\titlefont=cmr12 at 13pt\titlefont $\bullet$ 11:07: initial work}\smallskip
This one took a bit more time because of the roots, but it was ultimately a variation on one of the examples in 2.5 involving roots.
\par

Thm 2.7 I assume still holds, and the limit at infinity ends up being ${\root 3 \of x^6 \over 3x^6}$. Simplifying, ${\root 3 \of x^6} = (x^6)^{1/3}=x^2$, which yields ${x^2 \over 3x^2} = {1 \over 3}$.
\par

{\smallskip\noindent\bf\font\titlefont=cmr12 at 13pt\titlefont $\bullet$ 14:49: My limit was incorrect}\smallskip
It did not like ${1 \over 3}$.
\par

{\smallskip\noindent\bf\font\titlefont=cmr12 at 13pt\titlefont $\bullet$ 14:55: attempting this problem again}\smallskip
I forgot to factor in the other root in the denominator. Dividing by $x^2$ in the denominator is effectively dividing by $x^4$ under the root. So, the 5 remains, and you (possibly) get $1 \over {3 + \sqrt(5)}$. I'm hesitant on this, and will think on this before submitting.
\par

Still need to check the limits of this HA though.
\par

\xrdef{node16}
{\medskip\noindent\bf\font\titlefont=cmss17 at 20pt\titlefont 2 5 73}\medskip
Consider the function
$$\
f(x) = {
3x^4 + 12x^3 - 36x^2
\over
x^4 - 40x^2 + 144
}
$$
\par a. evaluate $$\lim_{x \to \infty}{f(x)}$$ and
$$\lim_{x \to -\infty}{f(x)}$$ and then identify the horizontal
asymptotes.
\par b. Find the vertical asymptotes. For each asymptote $x = a$,
evalulate $$\lim_{x \to a^{-}}{f(x)}$$ and $$\lim_{x \to a^{+}}{f(x)}$$
{\medskip\noindent\bf\font\titlefont=cmitt12 at 17pt\titlefont 2026-02-10}\medskip
{\smallskip\noindent\bf\font\titlefont=cmr12 at 13pt\titlefont $\bullet$ 12:01: initial thoughts}\smallskip
Finding limits at infinitiess again, and then identifying asymptotes.
\par

powers are the same, I hope that makes the limits easier (something something 3?)
\par

VAs:  looks factorable. Lots of squares.
\par

{\medskip\noindent\bf\font\titlefont=cmitt12 at 17pt\titlefont 2026-02-11}\medskip
{\smallskip\noindent\bf\font\titlefont=cmr12 at 13pt\titlefont $\bullet$ 11:22: initial work. swing and a miss}\smallskip
horizontal asymptotes are easy enough to do here. Getting the vertical asympotes are proving to be a pain for me. I have been trying to figure out ways to reduce the denominator.
\par

There are a lot of squares in here, which maybe implies some square root operations, but that leads to an ugly middle term. I left off trying to divide both polynomials by $x^2$ so they both end up having $x^2$ as their highest order polynomial. Again, there's an ugly term in the denom.
\par

Going to have to swing back for this one.
\par

{\smallskip\noindent\bf\font\titlefont=cmr12 at 13pt\titlefont $\bullet$ 13:33: more stabs on this}\smallskip
I tried a few approaches to reduce this damn denominator. Dividing everything by $x^2$ seems to get 144 into a term $(12/x)^2$, which seems awfully similar to the $12x$ term in the numerator.
\par

Frustrated. I didn't log my time, but it was at least half an hour, probably more. I hate finding the "trick".
\par

\xrdef{node17}
{\medskip\noindent\bf\font\titlefont=cmss17 at 20pt\titlefont 2 5 75}\medskip
Consider
$$\
f(x) = {
x^2 - 64
\over
x(x - 8)
}
$$
\par a. evaluate $$\lim_{x \to \infty}{f(x)}$$ and
$$\lim_{x \to -\infty}{f(x)}$$, and then identify the horizontal
asymptotes.
\par b. Find the vertical asymptotes. For each asymptote $x = a$,
evalulate $$\lim_{x \to a^{-}}{f(x)}$$ and $$\lim_{x \to a^{+}}{f(x)}$$
{\medskip\noindent\bf\font\titlefont=cmitt12 at 17pt\titlefont 2026-02-10}\medskip
{\smallskip\noindent\bf\font\titlefont=cmr12 at 13pt\titlefont $\bullet$ 12:06: initial thoughts}\smallskip
limits at inf, and asymptotes. looks factorable.
\par

There's some canceling that can be done, so I think the asymptote would be at $x = 0$.
\par

{\medskip\noindent\bf\font\titlefont=cmitt12 at 17pt\titlefont 2026-02-11}\medskip
{\smallskip\noindent\bf\font\titlefont=cmr12 at 13pt\titlefont $\bullet$ 13:36: initial work}\smallskip
First, I simplified and reduced the equation via factoring. Assuming 2.7 theorem still applies, the answer I think is just 1 at both sides? In the reduced form, the only VA is $x = 0$, because $x$ is the only thing in the denominator still.
\par

{\smallskip\noindent\bf\font\titlefont=cmr12 at 13pt\titlefont $\bullet$ 14:52: It did not like my vertical asymptotes}\smallskip
Similar to 2.5.78, it did not like that I put in 0. It should have been 1, not 0.
\par

{\smallskip\noindent\bf\font\titlefont=cmr12 at 13pt\titlefont $\bullet$ 15:09: repeating this problem again}\smallskip
I got the HA before, just need to get the VA correct. I think this was a small mistake.
\par

I've redone this, I've double checked the plots. It looks like the vertical asymptote was at 0, which was my initial. Maybe I got the limits backwards? RHS limit is positive, LHS is negatve.
\par

{\smallskip\noindent\bf\font\titlefont=cmr12 at 13pt\titlefont $\bullet$ 15:26: Looking at the answers, not sure what I am missing here}\smallskip
The plot definitely has one vertical asymptote at x=0, and there are only 2 options for this. One says the LHS limit doesn't exist, which it does.
\par

Either I am missing something, or I am putting in the answer wrong. grr.
\par

\xrdef{node18}
{\medskip\noindent\bf\font\titlefont=cmss17 at 20pt\titlefont 2 5 78}\medskip
Consider the function
$$\
f(x) = {
|98 - 2x^2|
\over
x(x + 7)
}
$$
\par a. Analyze $$\lim_{x \to \infty}{f(x)}$$ and
$$\lim_{x \to -\infty}{f(x)}$$, and then identify the horizontal
asymptotes.
\par b. Find the vertical asymptotes. For each asymptote $x = a$,
evalulate $$\lim_{x \to a^{-}}{f(x)}$$ and $$\lim_{x \to a^{+}}{f(x)}$$
{\medskip\noindent\bf\font\titlefont=cmitt12 at 17pt\titlefont 2026-02-10}\medskip
{\smallskip\noindent\bf\font\titlefont=cmr12 at 13pt\titlefont $\bullet$ 12:07: initial thoughts}\smallskip
Limits and horizontal asymptotes. Vertical asymptotes. Absolute value in numerator. Weird number 98. Wait, you can bring out a 2 and it becomes a 49, and there is an $(x + 7)$ in the denominator. So that's promising.
\par

{\medskip\noindent\bf\font\titlefont=cmitt12 at 17pt\titlefont 2026-02-11}\medskip
{\smallskip\noindent\bf\font\titlefont=cmr12 at 13pt\titlefont $\bullet$ 13:56: initial work}\smallskip
This has got an absolute value in the numerator. Examining this, I figured that the inside expression would always be negative as $x$ approaches infinity in both directions, so the whole expression becomess negated.
\par

Distributing values, the expression seems to want to have a horizontal asymptote at $y=2$.
\par

To get the vertical asymptotes, factor the expression in the numerator, and then cancel. There is only one VA left, and that is $x=0$. Finding the one sided limits, one notes that only negative values approach the denominator from the left side, and positive values only approach the numerator from the right hand side. So you get $-\infty$ and $+\infty$, respectively.
\par

{\smallskip\noindent\bf\font\titlefont=cmr12 at 13pt\titlefont $\bullet$ 14:50: It did not like my vertical asymptotes}\smallskip
I assumed there'd be a VA at 0 because I was dividing by x. I realized I didn't distribute right. I think it should be 2.
\par

{\smallskip\noindent\bf\font\titlefont=cmr12 at 13pt\titlefont $\bullet$ 15:30: re-doing 2.5.78}\smallskip
I plotted this. While there's definitely a vertical asymptote at $x = 0$, there's also something funky at $x=-7$. It doesn't look asympote-y though, more of a discontinuity. I'm not sure why. 
\par

\xrdef{node19}
{\medskip\noindent\bf\font\titlefont=cmss17 at 20pt\titlefont 2 6 15}\medskip
What is the domain of
$$\
f(x) = {
x^2
\over
x - 2
}
$$
and where is $f$ continuous?
{\medskip\noindent\bf\font\titlefont=cmitt12 at 17pt\titlefont 2026-02-11}\medskip
{\smallskip\noindent\bf\font\titlefont=cmr12 at 13pt\titlefont $\bullet$ 16:16: initial thoughts}\smallskip
This is asking for the domain where f is continuous for an equation that is a rational polynomial. The only time the function is undefined is when $x=2$. So I think it's continuous for all values when $x \ne 2$.
\par

\xrdef{node20}
{\medskip\noindent\bf\font\titlefont=cmss17 at 20pt\titlefont 2 6 17}\medskip
Determine whether the following function is continuous at
$a$. Use the continuity checklist to justify your answer.
$$\
f(x) = {
2x^2 + 3x + 1
\over
x^2 - 8x
}, a = 8
$$
{\medskip\noindent\bf\font\titlefont=cmitt12 at 17pt\titlefont 2026-02-11}\medskip
{\smallskip\noindent\bf\font\titlefont=cmr12 at 13pt\titlefont $\bullet$ 16:17: initial thoughts}\smallskip
Doesn't look like this function is going to be continuous at $a = 8$ because $f(8)$ seems to be undefined.
\par

\xrdef{node21}
{\medskip\noindent\bf\font\titlefont=cmss17 at 20pt\titlefont 2 6 21}\medskip
Determine whether the following function is continuous at $a$.
Use the continuity checklist to justify your answer.
$$
f(x) = \cases {
{x^2 - 9 \over x - 3} & if $x \ne 3$ \cr
{4} & if $x = 3$ \cr
} ; a = 3
$$
{\medskip\noindent\bf\font\titlefont=cmitt12 at 17pt\titlefont 2026-02-11}\medskip
{\smallskip\noindent\bf\font\titlefont=cmr12 at 13pt\titlefont $\bullet$ 16:19: initial thoughts}\smallskip
$f(3)$ is defined. Eyeballing, looks like the limits might not work out to be $f(a)$. LHS limit looks to be 6? RHS is 4.
\par

\xrdef{node22}
{\medskip\noindent\bf\font\titlefont=cmss17 at 20pt\titlefont 2 6 23}\medskip
Determine whether the following function is continuous at $a$.
Use the continuity checklist to justify your answer.
$$\
f(x) = {
5x - 2
\over
x^2 - 7x + 10
}, a = 2
$$
{\medskip\noindent\bf\font\titlefont=cmitt12 at 17pt\titlefont 2026-02-11}\medskip
{\smallskip\noindent\bf\font\titlefont=cmr12 at 13pt\titlefont $\bullet$ 16:21: initial thoughts}\smallskip
$f(a)$ doesn't seem to be defined at $a = 2$. So, this function is not continuous.
\par

\xrdef{node23}
{\medskip\noindent\bf\font\titlefont=cmss17 at 20pt\titlefont 2 6 25}\medskip
Determine the interval(s) on which the following function
is continuous.
$$
p(x) = 5x^4 + 6x^3 + 8
$$
{\medskip\noindent\bf\font\titlefont=cmitt12 at 17pt\titlefont 2026-02-11}\medskip
{\smallskip\noindent\bf\font\titlefont=cmr12 at 13pt\titlefont $\bullet$ 16:23: initial thoughts}\smallskip
This is just a polynonial expression, which is continuous. So the interval is going to be $(-\infty, \infty)$.
\par

\xrdef{node24}
{\medskip\noindent\bf\font\titlefont=cmss17 at 20pt\titlefont 2 6 28}\medskip
Determine the interval(s) on which the following function
is continuous.
$$
f(x) = {
x^2 - 6x + 5
\over
x^2 - 1
}
$$
Determine the interval(s) on which the following function
is continuous:
$$
f(x) = {
x^2 - 6x + 5
\over
x^2 - 1
}
$$
{\medskip\noindent\bf\font\titlefont=cmitt12 at 17pt\titlefont 2026-02-11}\medskip
{\smallskip\noindent\bf\font\titlefont=cmr12 at 13pt\titlefont $\bullet$ 16:24: initial thoughts}\smallskip
The denominator is $x^2 - 1$, so the function is not defined for $x = 1$ and $x = -1$. So, the intervals are from $(-\infty, -1)$, $(-1, 1)$, and $(1, \infty)$. In others continus when $x \ne -1$, $x \ne 1$.
\par

\xrdef{node25}
{\medskip\noindent\bf\font\titlefont=cmss17 at 20pt\titlefont 2 6 39}\medskip
Suppose $f(x)$ is defined as shown below.
\par a. Use the continuity checklist to show $f$ is not
continuous at 1.
\par b. Is $f$ continuous from the left or right at 1?
\par c. State the intervals of continuity.
$$\
f(x) = \cases{
x^2 + 3x & if $x > 1$ \cr
2x & if $x \le 1$ \cr
}
$$
{\medskip\noindent\bf\font\titlefont=cmitt12 at 17pt\titlefont 2026-02-11}\medskip
{\smallskip\noindent\bf\font\titlefont=cmr12 at 13pt\titlefont $\bullet$ 16:27: initial thoughts}\smallskip
a. function is defined at $x = 1$, but the limit is not defined.
\par

b. The piecewise function has  $x \le 1$, which I think means it is left continuous (but I know there's a formal definition and all that).
\par

c. intervals of continuity: $(-\infty, 1]$ and $(1, \infty)$.
\par

\xrdef{node26}
{\medskip\noindent\bf\font\titlefont=cmss17 at 20pt\titlefont 2 6 43}\medskip
Determine the interval(s) on which the following function
is continuous. At which finite endpoints of the intervals
of continuity is $f$ continuous from the left or continuous
from the right?
$$f(x) = \sqrt{3x^2 - 36}$$
{\medskip\noindent\bf\font\titlefont=cmitt12 at 17pt\titlefont 2026-02-11}\medskip
{\smallskip\noindent\bf\font\titlefont=cmr12 at 13pt\titlefont $\bullet$ 16:32: initial thoughts}\smallskip
Square root. At some point, the square root becomes negative, which doesn't exist for real numbers. Where that exists seems to be a non-integerer value. Somewhere between 3 and 4 (positive or negative). This might take some alegebra to figure out the endpoints precisely.
\par

\xrdef{node27}
{\medskip\noindent\bf\font\titlefont=cmss17 at 20pt\titlefont 2 6 61}\medskip
Determine the interval(s) on which the function
$f(x) = \csc x$
is continuous, then analyze the limits
$$\
\lim_{x \to -{\pi \over 3}}{f(x)}
$$
and
$$\
\lim_{x \to -{2\pi}^{-}}{f(x)}
$$
{\medskip\noindent\bf\font\titlefont=cmitt12 at 17pt\titlefont 2026-02-11}\medskip
{\smallskip\noindent\bf\font\titlefont=cmr12 at 13pt\titlefont $\bullet$ 16:35: initial thoughts}\smallskip
This is a trig one. $\csc = 1/\sin$. Why does this problem want you to find intervals and then limits. Looks confusing.
\par

\xrdef{node28}
{\medskip\noindent\bf\font\titlefont=cmss17 at 20pt\titlefont 2 6 65}\medskip
Determine the intervals on which the following function
is continuous. Then analyze the given limits.
$$\
f(x) = {
5e^x
\over
1 - e^{4x}
};
\lim_{x \to 0^{-}}{f(x)};
\lim_{x \to 0^{+}}{f(x)}
$$
{\medskip\noindent\bf\font\titlefont=cmitt12 at 17pt\titlefont 2026-02-11}\medskip
{\smallskip\noindent\bf\font\titlefont=cmr12 at 13pt\titlefont $\bullet$ 16:37: initial thoughts}\smallskip
Find intervals, then limits. Again, not sure what the logic is here for doing intervals then limits. Is it easiererr to find the limits. These are more transcendental functions.
\par

\xrdef{node29}
{\medskip\noindent\bf\font\titlefont=cmss17 at 20pt\titlefont 2 6 86}\medskip
Determine the value of the constant a for which the function
$f(x)$ is continuous at $-4$.
$$\
f(x) = \cases{
{x^2 + 6x + 8 \over x + 4} & if $x \ne -4$ \cr
a & if $x = -4$ \cr
}
$$
\bye
{\medskip\noindent\bf\font\titlefont=cmitt12 at 17pt\titlefont 2026-02-11}\medskip
{\smallskip\noindent\bf\font\titlefont=cmr12 at 13pt\titlefont $\bullet$ 16:38: initial thoughts}\smallskip
Ah okay, we did something similar to this in the lecture. The proper procedure involves using the continuity checklist. You end up evaluating the one sided limits at $x \to a$, which yields two equations in terms of $a$. set those equal, and solve for $a$.
\par

\bye
